% Original pages: 317--321
% Figures: Figure 6 placeholder; Figures 7--8 placeholders on original p. 321.
% Uncertainty log: VERIFY on original p. 317 for the top entry of the lepton doublet; final sentence on original p. 321 truncates at the page boundary.

handed quarks and leptons are singlets. So bare masses are impossible,
and the quarks and leptons get masses only from $SU(2) \times U(1)$
breaking. We do not know why $\langle \phi \rangle$ is so small (relative to the large
masses of physics) but if this could be explained, the lightness of
the quarks and leptons would follow.

The problem with the Higgs boson is that its bare mass does not
violate $SU(2) \times U(1)$, or any other gauge symmetry. Given any Higgs
multiplet $\phi^i$, the mass term $\sum_i \phi_i^* \phi^i$ is always gauge invariant. We
would like a \underline{symmetry} violated by $m_\phi$, to explain why $\phi$ is so light.
The symmetry will hopefully be spontaneously broken, but only on a
very small mass scale, to explain why $\langle \phi \rangle$ is tiny but non-zero.

% VERIFY: the top entry of this lepton doublet is slightly unclear in the scan; it appears to be \nu_e.
Supersymmetry can do this if it relates the Higgs doublet $\binom{\phi^0}{\phi^-}$
to fermions like $\binom{\nu_e}{e^-}_L$ whose masses violate $SU(2) \times U(1)$. Then
$m_\phi = 0$ as long as supersymmetry and $SU(2) \times U(1)$ are unbroken. And
$m_\phi$ is small, solving the old problem of the ``large numbers,'' if we
can understand that supersymmetry and $SU(2) \times U(1)$ are \underline{weakly} broken.

So that we can discuss these matters in a more tangible way,
we must have available the explicit form of some supersymmetric Lagrangians.
Let us consider first supersymmetric theories with fields
of spin zero and spin one half only. In such theories, one may have
an arbitrary number of complex scalar fields $A^i$, their supersymmetric
partners being left-handed spinor fields $\psi_L^i$. We introduce a function
$W$ that depends only on the $A^i$, not on their complex conjugates $A^*_j$ ---
in other words, $W$ is an analytic function of the $A^i$. $W$ is usually
called the ``superspace potential''.

The Lagrangian is $\mathcal{L} = \mathcal{L}_{\mathrm{kinetic}} + \mathcal{L}_{\mathrm{scalar}} + \mathcal{L}_{\mathrm{Yukawa}}$, where
\begin{equation}
\begin{aligned}
\mathcal{L}_{\mathrm{kinetic}} &= \partial_\mu A_i^* \partial^\mu A^i + \bar\psi_i\, i \not\!\partial\, \psi^i \\
\mathcal{L}_{\mathrm{scalar}} &= - \sum_i \left|\frac{\partial W}{\partial A^i}\right|^2 \\
\mathcal{L}_{\mathrm{Yukawa}} &= - \left(\frac{\partial^2 W}{\partial A^i \partial A^j}\, \psi^i_\alpha \psi^{\alpha j} + h.c.\right)
\end{aligned}
\tag{19}
\end{equation}

For how this construction was discovered, I refer you to the literature.\textsuperscript{2}
The easiest way to show that this describes a supersymmetric
theory is to demonstrate that the supersymmetry current $S^\mu$ is conserved.
This current is
\begin{equation}
S^\mu = \partial_\alpha A_i^* \gamma^\alpha \gamma^\mu \psi^i + i\, \frac{\partial W^*(A^*)}{\partial A_i^*}\, \gamma^\mu \psi_i^*
\tag{20}
\end{equation}
The first term in (20) we have already seen in free field theory;
the second term is added due to the interactions.

Note that, for a renormalizable theory, $W$ should be at most a
cubic function of the $A^i$. If $W$ is cubic, then (19) contains only
terms of dimension four or less, and this corresponds to a renormalizable
theory.

For our purposes, the most important part of (19) is the formula
for the scalar potential,
\begin{equation}
V(A^i, A_j^*) = \sum_i \left|\frac{\partial W}{\partial A^i}\right|^2
\tag{21}
\end{equation}
Let us ask, under what conditions is supersymmetry spontaneously
broken at the tree level? Evidently, if for some value of the $A^i$
the equations
\begin{equation}
\frac{\partial W}{\partial A^i} = 0
\tag{22}
\end{equation}
are simultaneously satisfied, then for this value of the fields the
potential energy vanishes, classically. Supersymmetry is unbroken.
On the other hand, if the equations (22) are inconsistent --- if they
are not satisfied for any choice of the $A^j$ --- then the minimum of the
potential is strictly positive, and supersymmetry is spontaneously
broken.

The simplest example is the O'Raifeartaigh model.\textsuperscript{9} There are
three fields, $A$, $X$, and $Y$, and the superspace potential is
\begin{equation}
W(A, X, Y) = g\, A Y + \lambda\, X (A^2 - M^2)
\tag{23}
\end{equation}

Here $g$, $\lambda$, and $M$ are constants. This theory is technically natural,
because of global symmetries.\begingroup\renewcommand{\thefootnote}{\fnsymbol{footnote}}\footnotemark[1]\endgroup\ The scalar potential is
\begin{equation}
\begin{aligned}
V(A, X, Y) &= \left|\frac{\partial W}{\partial A}\right|^2 + \left|\frac{\partial W}{\partial X}\right|^2 + \left|\frac{\partial W}{\partial Y}\right|^2 \\
&= g^2 |A|^2 + \lambda^2 |A^2 - M^2|^2 + |gY + 2 \lambda A X|^2
\end{aligned}
\tag{24}
\end{equation}
This potential is strictly positive, and supersymmetry is spontaneously
broken. In fact, $\partial W/\partial Y = 0$ only if $A = 0$, and $\partial W/\partial X = 0$
only if $A = \pm M$; these requirements are clearly inconsistent. If
$g/\lambda M$ is large enough, the minimum of the potential is at $A = 0$ and
the vacuum energy is $\lambda^2 M^4$, plus quantum corrections. Expanding
around the minimum of the potential (and making use of our previous
formulas for the Yukawa couplings as well as the scalar interactions),
it is easy to see that the bosons and fermions have unequal masses,
which is expected, since the positive ground state energy indicates
spontaneous breaking of supersymmetry.

We will have more to say about this model later, but for the
moment let us consider a model of another kind. Let us consider a
theory with $SU(5)$ symmetry, the only scalar field being a complex
field $A^i_{\ j}$ in the adjoint representation of $SU(5)$ (so $\operatorname{Tr} A = 0$). The
most general choice of $W$ would be
\begin{equation}
W = \frac{g}{3}\, \operatorname{Tr} A^3 + \frac{M}{2}\, \operatorname{Tr} A^2
\tag{25}
\end{equation}
where $g$ and $M$ are constants. The equations $\partial W/\partial A^i_{\ j} = 0$ give
\begin{equation}
g\left((A^2)^i_{\ j} - \frac{1}{5}\, \delta^i_{\ j}\, \operatorname{Tr} A^2\right) + M\, A^i_{\ j} = 0
\tag{26}
\end{equation}

\begingroup
\renewcommand{\thefootnote}{\fnsymbol{footnote}}
\footnotetext[1]{The theory is invariant under $A \to -A$, $Y \to -Y$ and under $Y \to e^{i\alpha}Y$, $X \to e^{-i\alpha}X$. It should be noted that any transformation under which $W$ changes only by an overall phase is a symmetry operation (the phase of $W$ cancels out of the scalar potential and can be removed from the Yukawa couplings by a chiral transformation).}
\endgroup

If we assume that $A$ can be diagonalized by a unitary transformation,\begingroup\renewcommand{\thefootnote}{\fnsymbol{footnote}}\footnotemark[1]\endgroup\ then it is easy to see that there are three solutions:
\begin{equation}
\begin{aligned}
A^i_{\ j} &= 0 \\
A^i_{\ j} &= \frac{M}{3g}
\left(
\begin{array}{r}
1 \\
1 \\
1 \\
1 \\
-4
\end{array}
\right) \\
A^i_{\ j} &= \frac{M}{g}
\left(
\begin{array}{r}
2 \\
2 \\
2 \\
-3 \\
-3
\end{array}
\right)
\end{aligned}
\tag{27}
\end{equation}
These three solutions correspond to the unbroken gauge groups $SU(5)$,
$SU(4) \times U(1)$, and $SU(3) \times SU(2) \times U(1)$, respectively. They each
correspond to unbroken supersymmetry, and they are exactly degenerate
at zero energy at least in this approximation (figure (6)), because
they were all found by requiring $\partial W/\partial A^i_{\ j} = 0$.

\begingroup
\renewcommand{\thefootnote}{\fnsymbol{footnote}}
\footnotetext[1]{It is really when $SU(5)$ is made into a gauge symmetry that this assumption is justified, because of extra terms that are then present in the scalar potential. Since $A$ is complex and not hermitian, it cannot necessarily be diagonalized by a unitary transformation.}
\endgroup

% TODO FIGURE: Original Fig. 6 appears here on p. 321.
\begin{center}
\fbox{\parbox[c][1.5in][c]{0.34\textwidth}{\centering TODO\\Figure 6 placeholder\\Original p.~321}}
\par\smallskip
Figure 6
\end{center}

Now, what really is the physics of this theory? It depends
entirely on the nature of the quantum corrections to the effective
potential. If really $E = 0$ for each of the three vacuum states, as
appears to be the case in the classical approximation, then this one
theory describes three different, inequivalent worlds. In one world,
the strong gauge group is $SU(3)$ and a baryon is made from three
quarks; in one world, the strong gauge group is $SU(4)$ and baryons are
bosons, made from four quarks; in one world the strong gauge group is
$SU(5)$ and a baryon is made from five quarks.

If the quantum corrections break supersymmetry in, say, two of
the three worlds (figure (7)), then the true vacuum is the one in which
$E = 0$ and the supersymmetry is \underline{not}
spontaneously broken.

If supersymmetry is spontaneously broken, and $E \neq 0$, in each
of the three worlds, then (figure (8)) the true vacuum is the one in
which supersymmetry is broken most weakly and the vacuum energy is
least.

% TODO FIGURES: Original Figs. 7 and 8 appear here on p. 321.
\begin{center}
\begin{minipage}{0.42\textwidth}
\centering
\fbox{\parbox[c][1.2in][c]{0.8\linewidth}{\centering TODO\\Figure 7 placeholder\\Original p.~321}}
\par\smallskip
Figure 7
\end{minipage}\hfill
\begin{minipage}{0.42\textwidth}
\centering
\fbox{\parbox[c][1.2in][c]{0.8\linewidth}{\centering TODO\\Figure 8 placeholder\\Original p.~321}}
\par\smallskip
Figure 8
\end{minipage}
\end{center}

One might usually guess that such a degeneracy would be resolved,
in perturbation theory, by loop diagrams. Perhaps the vacuum energy
is of order $\alpha$, coming from a one loop diagram, or of order $\alpha^2$, coming
from a two loop diagram. The loop diagrams would then be inducing
supersymmetry breaking. Some particles whose masses violate super-

% VERIFY: the final sentence is truncated at the original p. 321 page boundary and continues on p. 322, which is outside this assignment.
